\hwhat{H25's equal-time talk dial measures feedback within one minute. Coupling acts at the recipient's next model call, so HH262 asked whether a dial built on that read-out call shows larger gains that are still subcritical. Round 1 covered 66 non-reserved units in 35 periods, with 1.47M receiving calls. The read-out loop gain $g_{lag}$ counts the extra talk messages that one message causes through reading. Every unit is subcritical (all upper bounds $<0.40$). The gain is 0.13 in regime III and $\approx0$ elsewhere, and \textbf{the HH claim (lagged $>$ equal-time) fails} (7/35 periods). Round 2 located H50's factor 2. In regime III it comes from H50's pair-RD design; in regime I it is call-class composition. H111's fluctuations give an independent total gain of 0.149 [0.123, 0.175], which matches $g_{lag}$. Regime I shows no gain on its chat clock either.}

\hmodel{Linear response on the call clock (models 01, 09). Agent $i$'s call $c$ is a talk call ($Y_{ic}=1$) or not. $R^m_{ic}$ counts peer messages read at $c$ that were posted within the call's latency $d_c$ before its start. $P_{ic}$ counts messages posted within $d_c$ after the start (in flight, unreadable). $a_{i,d,\kappa}$ is an agent $\times$ day $\times$ call-class field; $\bar r$ is the readers per message and $\bar m$ the messages per talk call. Round 2 adds the Fano-implied gain: H111's collective Fano ratio $\Phi$ inverted through its sum rule.}

\hmath{\vspace{-6pt}\begin{gather*}
P(Y_{ic}{=}1)=a_{i,d,\kappa}+\beta_m R^m_{ic}+\beta_o R^o_{ic}+\beta_P P_{ic}+\gamma\cdot E_{ic}+\rho\,Y_{i,c-1}\\
J_1^*=\beta_m-\beta_P,\qquad g_{lag}=\bar m\,\bar r\,J_1^*,\qquad T/T_c=1/g_{lag}\\
G_H=\bar m\,\bar r\textstyle\sum_{h\le H}\delta(h),\quad \delta(h)=\sum_{k<h}\big[\beta(R^m_{c-k})-\beta(P_{c-k})\big]\\
\Phi_{obs}(10\,\mathrm{min})=\Phi_{pred}(g_{Fano})\approx(1-g_{Fano})^{-2}
\end{gather*}\vspace{-8pt}}

\hobs{../figures/summary_obs_col.pdf}{Round 1. Each point is one goal period (random-effects pool of its units, 95\% CI). The x-axis is H25's equal-time talk dial on the same calls, the y-axis the read-out loop gain. Regime-I periods sit on $g_{lag}\approx0$ while $g_{eq}$ reads 0.05--0.40: a fast shared field, not read-out coupling. Regime-III periods lie near the diagonal at 0.07--0.23.}

\htelos{For a monitor: report $g_{lag}$, not the equal-time dial or a pair-RD jump, as distance to a talk cascade. One regime-III talk message causes about 0.13 further messages through reading (amplification $\times1.15$, $T/T_c\approx8$). The fluctuation bound caps the full gain, all hops included, at 0.175. Naming carries the coupling: a named message raises the next-call talk probability by 0.08, an unnamed one by 0.004. In regime I, talk is set by the chat-turn schedule, not by reading. No village configuration came near runaway, so this is a gauge, not a lever.}

\hbottom{Read-out talk coupling is small and subcritical ($g_{lag}\approx0.13$ in regime III, matching the Fano-implied 0.149 [0.123, 0.175]; $\approx0$ in regime I on either clock). H50's twice-larger value is an artifact of its estimator design.}

\hresults{\scriptsize\setlength{\tabcolsep}{2pt}\begin{tabular}{@{}p{0.34\columnwidth}p{0.48\columnwidth}p{0.15\columnwidth}@{}}
\textbf{Prediction (dated)} & \textbf{Observed} & \textbf{Verdict}\\\hline
P1 subcritical everywhere & max upper bound 0.39; all units $<0.5$ & supported\\
P2 $g_{lag}>g_{eq}$ in $\ge2/3$ & 7/35 periods; median ratio 0.09 & failed\\
P3 rank-consistent with H42 & $\rho=0.55$ ($p=0.0007$) & supported\\
P5 named $\ge5\times$ unnamed (III) & 0.079 vs 0.004 ($\times19$) & supported\\
P7 $J_1^*$ falls with $N$ & $\rho=-0.59$; $g_{lag}$ flat & supported\\\hline
\multicolumn{3}{@{}l}{\textit{Round 2 (2026-10-05)}}\\
R3-P1 reproduce H50 & L1/L0 0.94, $\rho$ 0.87 & supported\\
R3-P2 gap in cells + placebo (III) & gap in pair RD $\to$ call counts (share 1.0); placebo 0.016 & failed\\
R3-P3 first stage $\Delta R\approx1$ & 1.04 & supported\\
R3-P4 cells remove $\ge$70\% (I) & 69\% (85--95\% other readouts) & failed narrowly\\
R4 multi-hop $G_5$ & synthetic failed; pool 0.24 [$-$0.03, 0.51] & descriptive\\
R4-P4 Fano check & $g_{Fano}$ 0.149 [0.123, 0.175] vs $g_{lag}$ 0.127 & hop 1 closes\\
R1-P1 chat-clock gain (I) $\ge$0.05 & $-$0.022 [$-$0.052, 0.009] & failed\\
R1-P2 explains H111 excess & $r_F$ 1.35 unchanged & failed\\
\end{tabular}}

\hobsb{../figures/round2_col.pdf}{Round 2. Left: unit-median gain along the ladder from H50's estimator to H67's. Regime I collapses at the call-class cells; regime III drops only at the switch from pair RD to per-call counts. Middle: multi-hop kernel $\delta(h)$, regime-III pool (descriptive; the design failed its synthetic test). Right: regime-I units on the chat clock (points, 95\% CI) and round 1's hop-1 gain (crosses). Both sit at 0.}

\hcaveats{The multi-hop design is too noisy per unit (SD 0.4--1.5 in synthetic worlds) and cannot resolve hop 1 in real data. The full-gain bound rests on H111's linear sum rule. Start-time error attenuates the chat-clock gain to about 0.23 of the truth, so regime I is bounded at $<0.04$, not zero. Post hoc: without conditioning on call mode, regime I shows 0.035 [0.023, 0.046] above a call-skeleton null; 3/4 of the "read $\to$ chat-mode" effect is the chat-turn schedule. The unnamed part sits at the placebo's resolution. Predictions were written knowing H25, H42, H50 and H111.}

\hnext{Confirm $g_{lag}\approx g_{Fano}$ on the reserved regime-III units. Identify the pair-sampling term (0.18 vs 0.13). Test the chat-turn schedule as the regime-I fast field behind H111's excess.}
